\chapter{Background and Related Work}
\label{ch:background}

\section{The turn problem inside a conversation between two humans}
\label{sec:turntaking}

Turn-taking is the way two speakers alternate in holding the floor. For a machine, the
difficulty is to decide, at every short pause in the exchange, whether to start speaking
because the turn has ended, or to keep waiting because the other party has not finished.
That decision rests on a trade-off that Section~\ref{sec:tradeoff} states formally, and no
setting escapes both of its errors. 

In most dialogue systems the decision is taken from
silence alone: the system waits for a pause of fixed length and then treats the turn as
over. The organisation of turn-taking has been studied since \citet{sacks1974turntaking},
and the rule is simple and fast, but it conflates two very different situations, a
genuine end of turn and a pause in the middle of an utterance, because both present
themselves as silence. In medical telephony this confusion is both frequent and costly: a
patient spelling a name or dictating a number produces silences that are not ends of turn.

These are precisely the cases on which models that listen only to sound fail. 

Having no access to the meaning of the words, they cannot separate a finished sentence from a
spelling sequence or a number still in progress. A perceptual study on French shows that the two channels are not interchangeable
\citep{uro2024perceptive}. Inter-pause units drawn from the REPERE corpus were annotated
for terminality and presented to participants either as audio or as text, with the end of
the unit truncated by up to three prosodic words; participants then judged whether the floor
could be taken. The audio condition allowed a turn end to be anticipated at least three
words before it arrived, the text condition considerably less.

That result ranks the two channels, and ranks them in favour of audio, which is the baseline
this project sets out to improve. What it establishes is separability: the same boundary
judged from sound and from words produces different judgements. Whether the difference
combines into something better than either channel alone is an empirical question, and
Section~\ref{sec:fusionresults} answers it.

\section{Tools and technologies}
\label{sec:tools}

The system reuses several existing components, described briefly here.

Voice activity detection is handled by Silero VAD \citep{silero_vad}, an open-source
package with no accompanying academic publication, which labels each short audio frame as
speech or silence. It is the entry point of the pipeline: it triggers an evaluation at
every transition from speech to silence.

The acoustic end-of-turn estimate comes from Smart Turn v3, a model released by Pipecat
and Daily \citep{smartturn_v3}, likewise without an academic publication and documented
only by its model card and a technical blog post. It uses the encoder layers of Whisper
Tiny \citep{radford2023whisper} with a single linear classification head on top, for
roughly eight million parameters. From the recent audio it produces a probability that the
turn has ended, without reading any text. This project uses version 3.2 in its CPU export.
Versions 3.1 and 3.2 improved the handling of short utterances and robustness to noise,
and the exact version is stated because the reported figures depend on it.

The lexical estimate comes from CamemBERT \citep{martin2020camembert}, a French
BERT-style language model applied to the partial transcript. Unlike the acoustic signal,
it judges completeness from the meaning of the words. Continuous transcription is provided
by a commercial speech recognition service (Azure), which emits partial hypotheses as
speech unfolds.

Evaluation relies on eot-bench \citep{eotbench}, which its authors describe as the first
open collection of genuine human-to-agent turns with aligned audio and textual context,
covering fourteen languages of which French is one. Rather than fixing a single threshold,
the harness sweeps the policy jointly over threshold, action delay and timeout and reports
the frontier between latency and cut-ins. It is an industrial artefact rather than a
peer-reviewed publication, released under Apache-2.0 in June 2026 alongside the LiveKit Turn
Detector v1, three weeks before implementation began here.

One caveat qualifies Chapter~\ref{ch:results}. Its authors omit scalar ranking metrics such
as the area under the receiver operating characteristic curve (AUC) from their comparison
reports, treating them as diagnostic only. This dissertation reports per-span AUC as its
headline figure because the question here is whether a lexical signal improves the
underlying \emph{score}, which a policy sweep conflates with the timing applied on top of
it. Both views are reported throughout.

\section{Calibration and cost-sensitive decision}
\label{sec:calibration-lit}

Neither of the two bodies of method underpinning the decision layer is specific to
turn-taking.

The first is probability calibration. A classifier is calibrated when, among the cases it
assigns a probability of $0.8$, roughly eight in ten are positive. Modern neural networks
are systematically miscalibrated and typically overconfident;
\citet{guo2017calibration} attribute this to model capacity and to training choices such as
batch normalisation, and show that accuracy and calibration can move in opposite directions.
Calibration is therefore not implied by good discrimination: a model can rank correctly, and
so obtain a high AUC, while the values it emits are meaningless as frequencies. 

The standard remedies are post-hoc and monotonic. Platt scaling fits a sigmoid; isotonic regression
\citep{zadrozny2002isotonic} fits an arbitrary non-decreasing step function, more flexible
but hungrier for held-out data \citep{niculescu2005predicting}. Being monotonic, neither
changes the ranking, and so neither changes the AUC: they correct the values, not the
ordering. Calibration quality is measured by the expected calibration error, or ECE
\citep{naeini2015ece,guo2017calibration}, defined in Section~\ref{sec:calibdesign}.

The second is cost-sensitive decision making. Once a probability is trustworthy, where to
act is not a free parameter but a consequence of the relative cost of the two errors:
\citet{elkan2001cost} shows that under an asymmetric loss the operating point follows from
the cost ratio rather than being set by hand. The rule used here decides on the evidence
available at one instant. The richer alternative lets observations accumulate and weighs
acting now against paying to observe more; the canonical formulation is Wald's sequential
probability ratio test \citep{wald1945sequential}, which Chapter~\ref{ch:conclusions}
returns to as future work.

A third strand is relevant to the in-domain dataset rather than the decision. Weak
supervision replaces manual annotation with cheap, noisy labelling heuristics, an approach
systematised by \citet{ratner2017snorkel}. It buys
scale and, in a clinical setting, avoids a person having to listen to call content, at the
cost of label noise that bounds the signal any model can recover from the data. That bound explains the in-domain result of Section~\ref{sec:gold}.

\section{Related work on end-of-turn detection}
\label{sec:related}

Work on end-of-turn detection can be grouped by the signals it uses.

\emph{Purely acoustic} approaches decide from sound, prosody and silence, without reading
the words. They are fast and independent of transcription, but blind to meaning, and so
cut off utterances that are grammatically incomplete. Smart Turn, the baseline of this
project, belongs to this family. A more elaborate member is
\citet{zink2024predictive}, which trains a Conformer to transcribe before an utterance
finishes and reads the decoder's cross-attention at the end-of-sentence token to locate the
endpoint, so no separate language model is needed. It is evaluated offline on pre-segmented
utterances, however, and never inside a running dialogue system, so it is not shown that the
timing it saves survives streaming constraints. That gap shaped the present design.

\emph{Purely lexical} approaches read the transcribed text and judge whether the utterance
is complete. They supply the meaning the acoustic signal lacks, but they inherit a dependency that shapes
both the design of Chapter~\ref{ch:approach} and the failure reported in
Section~\ref{sec:gold}. A lexical branch is only as good as the transcript it reads, and a
streaming transcript is late, unstable and imperfect: partials arrive behind the audio, are
revised as context arrives, and degrade on narrowband telephone speech. Waiting for a stable
transcript buys accuracy with the latency endpointing is trying to save; reading whatever is
available reads a text that may be wrong. Neither horn is free.

\emph{Joint} approaches process audio and text together. FastTurn \citep{fastturn2026} unifies streaming connectionist temporal classification (CTC)
decoding with acoustic features so that a
decision can be taken early from partial observations while semantic cues are preserved.
It performs well but is heavy and presupposes joint training of its components. JAL-Turn
\citep{jalturn2026} combines pre-trained acoustic representations with linguistic features
through a lightweight joint model built on a frozen automatic speech recognition (ASR)
encoder, so that it runs in
parallel with recognition and adds no latency. Unlike FastTurn it does not require the two
modalities to be retrained together, and its labels are extracted automatically from
large-scale dialogue corpora. 

X2-Turn \citep{x2turn2026} goes further, placing a frame-synchronous turn state head
alongside the ASR head on shared streaming features. Phoenix-VAD \citep{phoenixvad2025}
takes the opposite route to the one adopted here, putting a large language model inside the
endpoint decision itself.

On the benchmarking side, the LiveKit detector and test set provide the policy-sweep
evaluation framework this project adopts. Other work remains purely acoustic: endpoint
anticipation \citep{udupa2026anticipation} bypasses ASR to predict a turn ending in
advance, up to roughly 2.5 seconds ahead, so that the rest of the pipeline can be started
sooner. Its object is not the choice of a decision threshold but early prediction, and the
cost it seeks to reduce is the speculative computation wasted when the pipeline is
launched too early and then abandoned.

The closest neighbour is Easy Turn \citep{easyturn2025}, an open-source modular detector
that also integrates acoustic and linguistic information and also positions itself against
Smart Turn. 

Beyond that the two diverge on every axis that matters. Easy Turn trains: its
model is fitted on a purpose-built 1{,}145-hour corpus and predicts one of four dialogue
states, from which a decision is taken by argmax. The present work trains no base model,
fitting only a linear head on frozen embeddings and a four-parameter fusion above it, and
emits a scalar probability rather than a class. 

Easy Turn reports accuracy on its own test set; this work reports on a third-party benchmark in a language subset it did not choose.
And neither Easy Turn nor any system reviewed here reports whether its probability is
calibrated, or connects its operating point to a stated cost. Fusion is done in both; what
differs is that the output is treated as a probability to be validated before a threshold
placed on it can mean anything.

Table~\ref{tab:related} summarises the comparison.

\begin{table}[H]
\centering
\caption[Recent turn-taking systems compared]{Representative recent end-of-turn detection
systems and where this work sits. ``Calib.'' records whether the system reports any
measurement or correction of the calibration of its own output probability.}
\label{tab:related}
\footnotesize
\setlength{\tabcolsep}{3pt}
\begin{tabular}{@{}p{2.2cm}p{2.4cm}p{2.5cm}p{2.8cm}p{1.0cm}p{2.5cm}@{}}
\toprule
System & Signals & Training & Decision rule & Calib. & Evaluation \\
\midrule
Easy Turn (2025) & acoustic + linguistic & trained, 1{,}145\,h & argmax over 4 states & no & own test set \\
Phoenix-VAD (2025) & LLM over text & LLM fine-tuned & LLM in the loop & no & own test set \\
JAL-Turn (2026) & acoustic + linguistic & head only, frozen ASR & fixed threshold & no & multilingual \\
FastTurn (2026) & acoustic + CTC & joint training & early, on partials & no & own test set \\
X2-Turn (2026) & shared ASR features & joint, dual-head & per-frame state & no & own test set \\
Endpoint antic. (2026) & acoustic only & trained & per-horizon threshold & no & anticipation only \\
LiveKit v1 (2026) & text (+ VAD) & trained detector & swept policy & no & 14-language benchmark \\
\midrule
\textbf{This work} & acoustic + lexical & \textbf{frozen bases}; linear head + fusion & \textbf{cost-derived, on a calibrated probability} & \textbf{yes} & French benchmark + medical domain \\
\bottomrule
\end{tabular}
\end{table}

The fusion mechanism itself, combining an acoustic signal with a lexical one, and the
latency-versus-cut-in framing are not new, and this project claims to have invented
neither. The field is moving fast enough that every one
of the seven systems in Table~\ref{tab:related} was published in the eleven months to
August 2026, and the most recent of them appeared three weeks before this chapter was
written.

\section{The gap addressed}
\label{sec:gap}

Three gaps emerge from this review. Almost all of the work above appeared after this
project was scoped (Chapter~\ref{ch:intro}), so these are the gaps that remain \emph{after}
that shift, not the ones the project was originally scoped against.

First, none of the systems reviewed reports an explicit
calibration of its probability, neither an expected calibration error nor any
recalibration, and none relates its decision threshold to a cost of interruption against
latency: eot-bench sweeps thresholds and some vendors expose confidence scores, but the
probability itself is neither measured nor corrected, so a chosen threshold has no precise
meaning. 

Second, Table~\ref{tab:related} shows the field moving steadily towards larger jointly
trained systems. The opposite end of the design space, where frozen components are merely
composed, is little studied.
Third, the French medical telephony domain, where semantically incomplete utterances
(spelling, digits, dates) are both frequent and consequential, is not addressed. These
are the three gaps this dissertation sets out to close.
